\section{The Inferior States of the Being}

\begin{quotex}
O put not your trust in princes, nor in any child of man, for there is no help in them.

For when the breath of man goeth forth, he shall turn again to his earth, \textbf{and then all his thoughts perish}. \flright{\textsc{Psalm 146:2-3)}}

\end{quotex}
In \textbf{Rene Guenon}`s metaphysical trilogy, the inferior states of the being are passed over without much comment. They are worth addressing here, as these will help us understand our interior life, as well as the Theological concept of Hell. Julius Evola spends some time describing these states in \emph{Revolt Against the Modern World:}

\begin{quotex}
According to esoteric teachings, at the death of the body an ordinary person usually \textbf{loses his or her personality}, which was an illusory thing even while that person was alive. The person is then reduced to a shadow that is itself destined to be dissolved after a more or less lengthy period culminating in what was called ``the second death." The essential vital principles of the deceased return to the totem, which is a \textbf{primordial, perennial, and inexhaustible matter}… the [infernal] path is that trodden by those who do not survive in a real way, and who slowly yet inexorably dissolve back into their original stocks… this is the life of Hades. \flright{\emph{Revolt Against the Modern World, Ch. VIII}}

\end{quotex}
These are the sub-individual states, spiraling ever downward to ``substance" or Chaos.  Here reside the Titans, the Cthonic gods of antiquity, and the vital forces of genus, species, and races.  Just as the supra-individual states of formless manifestation transcend individuality in the direction of luminous, all-comprehending intelligence, so too are the sub-individual states non-individual, but in the opposite direction: toward the blind ``forces" of the Waters, which make up the substance of all formal and non-formal manifestations.

To descend to the lower states means to reduce one's identity down to one of these sub-individual ``forces", e.g., to identify only as a member of a race, a class, or with one's individual proclivities, sexual obsessions, etc.  These are all forms of sub-human identity that dispense of the individuality, and entail real \textbf{identification }with such ``forces" below.

Upon identification with these states, ``reason" is lost, and the possibility of contact with the Personality is cut off.  Free will is thus made effectively impossible, and the human being becomes instead the plaything of the blind forces of the Waters, spiraling indefinitely downward until the ``second death" (Rev. 21:8).

\paragraph{King Lear}
Metaphysical teachings can seem dry and intellectual, unless we enliven them with imagination and experience.  Shakespeare's \emph{King Lear} can do just this service for us, if we take it as our guide through the lower states, akin to Dante's \emph{Inferno}.  The titular King Lear is a man who should be the supreme governor of his realm, but instead has lived poorly, and thus has gone mad in his old age.  He embarks on a downward spiral that will lead to the ruin of his kingdom.

\begin{quotex}

GONERIL: You see how full of changes his age is; the observation we have made of it hath not been little…

REGAN: 'Tis the infirmity of his age. Yet \textbf{he hath ever but slenderly known himself.}

GONERIL: \textbf{The best and soundest of his time hath been but rash}. Then must we look from his age to receive not alone the imperfections of long-engraffed condition, but therewithal the \textbf{unruly waywardness} that infirm and choleric years bring with them.

\flright{\emph{King Lear}, Act 1, Scene 1}

\end{quotex}
Because King Lear wasted his best years without ever striving to \textbf{know himself}, he has now arrived at old age with an ever-decreasing possibility of rising out of his low condition. As his mind and body deteriorate, his ability to sustain a link with his Personality has severely weakened, and so an ever more evident madness is the result. Recall Tomberg's teaching\footnote{\url{https://gornahoor.net/?p=13778}} that only ``moral memory" can maintain a sound mind into old age. Boris Mouravieff elaborates:

\begin{quotex}
Memory is a direct function of the \emph{being} of the individual… Loss of memory… makes a madman out of a normal man. 

\end{quotex}
King Lear becomes a shadow of himself. 

\begin{quotex}
GONERIL: I would you would make use of that good

Wisdom whereof I know you are fraught, and put away

\textbf{These dispositions that of late transform you}

\textbf{From what you rightly are.}

FOOL: \textbf{May not an ass know when the cart draws the horse?}

Whoop, Jug, I love thee!

LEAR: \textbf{Doth any here know me? This is not Lear.}

Doth Lear walk thus? speak thus? Where are his eyes?

Either his notion weakens, his discernings

Are lethargied Ha! waking? `Tis not so!

Who is it that can tell me who I am?

FOOL: \textbf{Lear's shadow.}

\flright{\emph{King Lear}, Act 1, Scene 4}
\end{quotex}

Falling into the lower states consists in being \textbf{transformed from what one rightly is}, hence the Fool's apt image of the \textbf{cart drawingt the horse}.  The Self, which should be active, has become passive to the forces of the lower states.  Ultimately, one does not even really \emph{have} a ``Self" anymore, and is left as a mere shadow or ``shade", which is to say, a mere psychic residue, which was only formerly the manifestation of a Person.

\begin{quotex}
GONERIL:

\textbf{As you are old and reverend, should be wise.}

Here do you keep a hundred knights and squires,

Men so disordered, so debauched and bold

That this our court, infected with their manners,

Shows like a riotous inn. Epicurism and lust

Make it more like a tavern or a brothel

Than a graced palace. The shame itself doth speak

For instant remedy.

…You strike my people, and \textbf{your disordered rabble}

\textbf{Make servants of their betters}. \flright{\emph{King Lear}, Act 1, Scene 4}

\end{quotex}
Lear should be wise at his age, but is not.  His court, which signifies also his inner life, is instead full of a huge, unruly retinue, like Mouravieff's mass of ``Little I's", which cause inner chaos and rule over him.  In Shakespeare's words, they ``make servants of their betters".  In other words, the ``retinue" enslaves the Individuality and the reasoning faculty.

Such is the nature of a man in thrall to the lower states of his being.  In life, it is always possible to emerge from these states by self-remembrance, penitence, and purification. It is even possible to dominate them, as Dante demonstrates. However, it is obvious that if one dies in identification with one of these states, escape from them \emph{postmortem} becomes quite impossible, because of the abandonment of Personhood and free will that such an identification entails.

\begin{quotex}
\textbf{To free human beings from the dominion of the totems}; to strengthen them; to address them to the fulfillment of a spiritual form and a limit; and to bring them in an invisible way to the line of influences capable of creating a destiny of \textbf{heroic and liberating immortality} — this was the task of the aristocratic cult. When human beings persevered in this cult, the fate of Hades was averted… Once the divine rites were neglected, however, this destiny was reconfirmed \textbf{and the power of the inferior nature became omnipotent again}… those who neglect the rites cannot escape ``hell," \textbf{this word meaning both a way of being in this life and a destiny in the next}. \flright{\emph{Revolt Against the Modern World, Ch. VIII}}

\end{quotex}
\textbf{Oremus et pro fratribus, vivis atque defunctis.}

\url{https://youtu.be/3mJdk5kOcds?t=446}



\flrightit{Posted on 2022-10-31 by Tannheuser }

\begin{center}* * *\end{center}

\begin{footnotesize}\begin{sffamily}



\texttt{Alex on 2022-11-05 at 19:44 said: }

Knowing oneself seems to involve solidifying a center of consciousness which is not easily shaken or undone. The modern life does not support the development of the center, the spiritual center which is the vehicle for the next stage of life. What we see with modern people is their reactionary body, emotions, feeble understanding of the depth of life. This reactionary body will dissolve with time and will be lost. Perhaps this is why most spirit communications deal with fragments of information, making anything meaningful difficult. Those living in the astral groves will not remain there very long as there is no point to it. So the recognition of a deep core of being will be that which moves us to higher states, not to be lost.


\end{sffamily}\end{footnotesize}
